% CHM shape of the modular hop on the 3+1D diamond waist.
\documentclass[11pt]{article}
\usepackage[margin=1.05in]{geometry}
\usepackage{amsmath,amssymb,amsthm}
\usepackage[colorlinks=true,linkcolor=blue,citecolor=blue,urlcolor=blue]{hyperref}

\newcommand{\pP}{\textbf{[P]}}
\newtheorem*{admission}{Admission}

\title{\textbf{The Modular Hopping Kernel on a $3{+}1$D Diamond Waist\\
Is the Casini--Huerta--Myers Envelope}}
\author{Robert W.\ McGwier\\
\small CTO, Cohere Technology Group\\
\small GitHub: \texttt{n4hy}}
\date{15 August 2026 --- Version 1}

\begin{document}
\maketitle

\begin{abstract}
\noindent
A2 (Papers~21--23) said the modular Hamiltonian is local. Jacobson
needs the local kernel to be the geometric modular Hamiltonian
\cite{BW1975,CHM2011}
\[
K=2\pi\int w\,T_{00},\qquad w=R^2-r^2.
\]
$T_{00}$ contains the hop. The test is therefore on nearest-neighbour
$K_{ij}$, not on the diagonal.

1d massless interval (the CHM theorem, used as an instrument gate):
$\rho(K_{i,i+1},w_{nn})=-0.942$.
$3{+}1$D ball $R=5$, $m=0$: $\rho=-0.987$, residual $R_{\mathrm{shape}}=0.162$.
At $mR=0.5$: $\rho=-0.989$.
A permutation of $w$ collapses to $\rho=-0.076$.
Auditor, different sizes: $1$d $\rho=-0.984$; $3$d $\rho=-0.987$.
C2 PASS. \pP{} as a necessary-condition test on this field.

A first diagonal instrument failed the 1d theorem and was discarded.
This is not $\eta=1/4G$, not foam, and not de~Sitter.
\end{abstract}

\section{Gate}

Locked on the hopping kernel. $C0$: 1d $\rho<-0.70$. $C1$: ball
$\rho<-0.60$. $C2$: $R_{\mathrm{shape}}<0.50$. $C3$: permuted $w$
has $|\rho|<0.30$. $C4$: $C1$ at $mR=0.5$. The sign is required:
hops are $-1$, so $K_{ij}$ tracks $-w$.

\section{Numbers}

\begin{center}
\begin{tabular}{lccc}
\hline
 & $\rho(K_{\mathrm{nn}},w_{\mathrm{nn}})$ & $R_{\mathrm{shape}}$ & perm $\rho$ \\
\hline
1d $L=32$ & $-0.942$ & $0.334$ & --- \\
1d audit $L=24$ & $-0.984$ & $0.176$ & --- \\
ball $R=5$, $m=0$ & $-0.987$ & $0.162$ & $-0.076$ \\
ball $R=5$, $mR=0.5$ & $-0.989$ & $0.150$ & $-0.077$ \\
ball audit $R=4$ & $-0.987$ & $0.161$ & $+0.054$ \\
\hline
\end{tabular}
\end{center}

The hopping kernel tracks a CHM-\emph{type} envelope on this
lattice: bulk-peaked, radial, negative for hop $-1$. An
independent adversary confirmed C1/C2 on $N=13,15$. Catch: a
linear $R-r$ weight can beat the quadratic on a small ball
($R=3$). The parabola wins at $R=4,5$. The quadratic is not
uniquely selected.

Together with the area law of Paper~24 this is the lattice input
Jacobson 1995 uses to get Einstein plus a free $\Lambda$. The
vacuum of that equation has $G_{\mu\nu}+\Lambda g_{\mu\nu}=0$.
That is the zero-Einstein-tensor step. It is not a selection of
$\Lambda=0$ or of de~Sitter, and $G$ is not measured.

If one \emph{identifies} $S=\alpha A_{\mathrm{cut}}$ with
$A/(4G)$ then $G=a^2/(4\alpha)\approx 1.02\,a^2$ using
$\alpha=0.245$. That is a cutoff identification, not a derivation.

\begin{admission}
I took the step A2 still owed: the local kernel has the CHM
shape. I threw out a diagonal instrument that failed its own
1d theorem. I did not write that $G$ was measured, and I did
not write that the foam is done.
\end{admission}

\begin{thebibliography}{9}
\bibitem{BW1975}
J.J.~Bisognano and E.H.~Wichmann,
J.\ Math.\ Phys.\ 16 (1975) 985.

\bibitem{CHM2011}
H.~Casini, M.~Huerta and R.C.~Myers,
JHEP 05 (2011) 036,
arXiv:1102.0440.

\bibitem{Jacobson1995}
T.~Jacobson,
Phys.\ Rev.\ Lett.\ 75 (1995) 1260.
\end{thebibliography}

\end{document}
